% =====================================================================
%  Chapitre 3 - Environnement de travail et architecture
% =====================================================================
\chapter{Environnement de travail et architecture}
\label{chap:architecture}

\section*{Introduction}
\addcontentsline{toc}{section}{Introduction}

Les exigences etant posees, ce chapitre presente les moyens mis en oeuvre
pour y repondre. Il decrit d'abord l'environnement de travail, materiel
et logiciel, puis l'architecture de la plateforme sous ses deux vues
complementaires : la vue logique, qui organise les responsabilites, et la
vue physique, qui decrit le deploiement. Il se termine par l'examen des
trois choix structurants qui caracterisent cette architecture, y compris
ceux qui s'ecartent du patron de reference.

% =====================================================================
\section{Environnement materiel}

Le developpement a ete conduit sur un poste unique, qui heberge
egalement l'ensemble des conteneurs de la plateforme.

\begin{table}[H]
\centering
\caption{Caracteristiques du poste de developpement}
\label{tab:materiel}
\small
\begin{tabularx}{\textwidth}{@{}p{3.4cm}X@{}}
\toprule
\textbf{Element} & \textbf{Caracteristique} \\ \midrule
Systeme d'exploitation & Windows 11 Professionnel \\
Environnement conteneurise & Docker Desktop, avec le sous-systeme Linux \\
Memoire vive & 16 Go, dont environ 8 Go alloues aux conteneurs \\
Stockage & Disque a semi-conducteurs, indispensable aux temps de
  compilation des huit services \\
\bottomrule
\end{tabularx}
\end{table}

Ce point merite d'etre souligne : \essentiel{une architecture
microservices se developpe difficilement sans conteneurisation}. Vingt
conteneurs cohabitent, dont une base de donnees, un courtier de messages
et une base vectorielle. Les installer et les configurer un a un sur le
poste serait long, fragile et non reproductible.

% =====================================================================
\section{Environnement logiciel}

\subsection{Outils de developpement et de gestion de version}

\begin{table}[H]
\centering
\caption{Outils utilises}
\label{tab:outils}
\small
\begin{tabularx}{\textwidth}{@{}p{3.2cm}p{3.0cm}X@{}}
\toprule
\textbf{Outil} & \textbf{Usage} & \textbf{Motif du choix} \\ \midrule
Visual Studio Code & Edition du code &
  Un seul editeur pour C\#, TypeScript, Python et SQL. \\
Docker Desktop & Execution locale &
  La definition des services est versionnee : le poste et le serveur
  executent la meme chose. \\
Git et GitHub & Gestion de version &
  Historique, branches et sauvegarde. \\
Postman & Essais d'API &
  Verification des points d'entree REST avant toute integration
  d'interface. \\
pgAdmin & Administration de la base &
  Inspection du schema et des donnees. \\
PlantUML & Modelisation UML &
  Les diagrammes sont du texte : ils se versionnent et se comparent
  comme du code. \\
Trello & Suivi du carnet &
  Voir la section~\ref{sec:outils-pilotage}. \\
\bottomrule
\end{tabularx}
\end{table}

\subsection{Langages, cadriciels et briques techniques}

\begin{table}[H]
\centering
\caption{Technologies retenues et versions}
\label{tab:technologies}
\small
\begin{tabularx}{\textwidth}{@{}p{2.9cm}p{2.6cm}X@{}}
\toprule
\textbf{Domaine} & \textbf{Briques et versions} & \textbf{Role} \\
\midrule
Services metier & .NET / C\# 8.0, ASP.NET Core 8.0, Entity Framework
  Core et Npgsql 8.0.11, MassTransit 8.3.6, SignalR, BCrypt.Net 4.0.3,
  FluentValidation 11.9.0, Serilog 8.0.3 &
  Socle des huit services metier, exposition REST, correspondance
  objet-relationnel, messagerie et journalisation \\
Interfaces & React 18.3 et Vite 6.0, Next.js 14.2, Expo 54 avec React
  Native 0.81, TypeScript 5.6, Leaflet 1.9, ApexCharts 4.4 &
  Back-office, portail client, application chauffeur, cartographie et
  graphiques \\
Aide a la decision & Python 3.11 et 3.12, FastAPI 0.115, XGBoost 2.1.3,
  scikit-learn 1.6.1, fastembed 0.4.2 &
  Estimation tarifaire, optimisation, prediction de pannes, assistant
  documentaire \\
Infrastructure & PostgreSQL 16, RabbitMQ 3.13, Qdrant 1.12, nginx 1.25
  et 1.27, Docker Compose &
  Persistance, courtier de messages, base vectorielle, passerelle et
  orchestration des vingt conteneurs \\
\bottomrule
\end{tabularx}
\end{table}

\begin{figure}[H]
  \centering
  \includegraphics[height=1.15cm,keepaspectratio]{logo-dotnet.png}\hspace{0.7cm}
  \includegraphics[height=1.15cm,keepaspectratio]{logo-csharp.png}\hspace{0.7cm}
  \includegraphics[height=1.15cm,keepaspectratio]{logo-react.png}\hspace{0.7cm}
  \includegraphics[height=1.15cm,keepaspectratio]{logo-typescript.png}\hspace{0.7cm}
  \includegraphics[height=1.15cm,keepaspectratio]{logo-python.jpg}\hspace{0.7cm}
  \includegraphics[height=1.15cm,keepaspectratio]{logo-postgres.png}\hspace{0.7cm}
  \includegraphics[height=1.15cm,keepaspectratio]{logo-rabbitmq.png}\hspace{0.7cm}
  \includegraphics[height=1.15cm,keepaspectratio]{logo-docker.png}
  \caption{Principales technologies mises en oeuvre}
  \label{fig:technologies}
\end{figure}

% =====================================================================
\section{Architecture logique}
\label{sec:archi-logique}

L'architecture logique organise les responsabilites en quatre niveaux
empiles, le flux descendant de l'interface vers la persistance.

\figReelle{logical-architecture.png}
  {Architecture logique de la plateforme, en quatre niveaux}
  {fig:archi-logique}{0.97}

\paragraph{Niveau 1, presentation.}
Trois applications servent trois publics : le back-office pour le
personnel, le portail pour les clients, l'application mobile pour les
chauffeurs. Aucune ne contient de regle metier ; elles consomment des API.

\paragraph{Niveau 2, passerelle.}
Un serveur nginx assure trois fonctions : le routage vers le service
concerne d'apres le prefixe de l'URL, la validation du jeton
d'authentification et le filtrage des origines autorisees. Il constitue
le point d'entree unique du systeme.

\paragraph{Niveau 3, coeur applicatif.}
Huit services metier en .NET et quatre services d'aide a la decision en
Python, relies par un courtier de messages. Une bibliotheque commune
porte les enumerations et les evenements d'integration partages.

\paragraph{Niveau 4, persistance.}
Une base relationnelle, une base vectorielle et un volume de fichiers.

\begin{definitionbox}{Une lecture a retenir}
Sur ce diagramme, les liaisons vont \exemple{de groupe a groupe} et non
de service a service. Les huit services metier partagent la meme
passerelle, le meme courtier et la meme base : les relier un a un
produirait une trentaine de traits sans rien apprendre de plus. Le detail
service par service figure dans les chapitres de sprint.
\end{definitionbox}

% =====================================================================
\section{Architecture physique}
\label{sec:archi-physique}

La vue physique decrit le deploiement effectif : vingt conteneurs sur un
hote unique, orchestres par Docker Compose.

\figReelle{physical-architecture.png}
  {Architecture physique : deploiement conteneurise}
  {fig:archi-physique}{0.97}

Trois points de cette vue meritent un commentaire.

\paragraph{La resolution des noms.}
A l'interieur du reseau, les services s'appellent par leur nom et non par
une adresse. Les ports 5001 a 5009 visibles sur le diagramme ne sont que
des publications vers l'hote, utiles au diagnostic et aux appels directs
du portail client.

\paragraph{L'ordre de demarrage.}
Il repose sur une condition de bonne sante : la base et le courtier
demarrent d'abord, puis les services metier, puis la passerelle. Les huit
services .NET exposent un point de controle de sante que
l'orchestrateur interroge.

\paragraph{Les acces directs.}
Le portail client et l'application chauffeur ne passent pas systematiquement
par la passerelle : le suivi de position et l'estimation tarifaire sont
atteints en direct. Ce n'est pas un oubli mais une consequence du
perimetre de routage de la passerelle, examinee en
section~\ref{sec:ecarts}.

% =====================================================================
\section{Choix structurants}
\label{sec:choix}

\subsection{Clean Architecture au sein de chaque service}

Chaque service metier est decoupe en quatre projets, dont les
dependances de compilation sont contraintes.

\figReelle{platform-clean-architecture.png}
  {Decoupage en couches d'un service metier}
  {fig:clean-archi}{0.52}

La regle est la suivante : \essentiel{les dependances pointent vers
l'interieur}. La couche \exemple{Domain}, qui porte les entites et les
regles de gestion, ne depend de rien d'autre que de la bibliotheque
commune. Elle ignore la base de donnees, le protocole HTTP et le
courtier de messages.

L'inversion se joue dans la couche \exemple{Infrastructure} : celle-ci
\essentiel{implemente} des interfaces declarees dans la couche
\exemple{Application}, et non l'inverse. Concretement, un service
applicatif manipule une interface de depot sans savoir qu'elle est
satisfaite par Entity Framework Core.

L'interet n'est pas theorique. Le domaine se teste sans base de donnees,
et le remplacement d'une brique technique n'affecte pas les regles
metier. Le service de maintenance predictive en fournit l'illustration :
son moteur de prediction est declare par une interface, et l'implantation
actuelle, fondee sur des regles deterministes, peut ceder la place a un
modele d'apprentissage sans qu'aucune autre couche ne bouge.

\subsection{Deux bus de messages plutot qu'un}

La plateforme utilise le meme courtier RabbitMQ de deux manieres
distinctes, ce qui surprend a la premiere lecture.

\begin{table}[H]
\centering
\caption{Les deux usages du courtier de messages}
\label{tab:deux-bus}
\small
\begin{tabularx}{\textwidth}{@{}p{3.2cm}XX@{}}
\toprule
 & \textbf{MassTransit} & \textbf{Echange par sujet} \\ \midrule
Format & Message type, avec enveloppe & JSON brut, sans enveloppe \\
Emetteurs & Services .NET & Services .NET \\
Destinataires & Services .NET & Services Python \\
Motif & Contrat fortement type entre services d'une meme pile &
  Les consommateurs Python liraient difficilement l'enveloppe
  MassTransit ; le JSON brut leur est directement exploitable \\
\bottomrule
\end{tabularx}
\end{table}

Ce choix est un compromis assume : il introduit deux conventions a
maintenir, mais evite d'imposer aux services Python le decodage d'un
format concu pour l'ecosysteme .NET.

\subsection{Ecarts assumes par rapport au patron de reference}
\label{sec:ecarts}

Une architecture microservices de reference prevoit une base de donnees
par service. \essentiel{Ce n'est pas le cas ici : les neuf chaines de
connexion pointent toutes vers la meme base.} Cet ecart doit etre
explicite plutot que masque.

\begin{table}[H]
\centering
\caption{Ecarts par rapport au patron de reference}
\label{tab:ecarts}
\small
\begin{tabularx}{\textwidth}{@{}p{3.6cm}X@{}}
\toprule
\textbf{Ecart} & \textbf{Portee et consequence} \\ \midrule
Base partagee &
  Le cloisonnement est LOGIQUE : chaque service ne possede en ecriture
  que ses propres tables et lit celles des autres domaines en SQL direct.
  C'est pourquoi certaines entites sont redeclarees en lecture seule dans
  plusieurs couches Domain. Consequence : le decouplage des donnees
  repose sur une discipline de developpement, non sur l'infrastructure. \\
Absence de cache &
  Aucun cache applicatif n'est deploye. La limitation de debit de
  l'assistant est tenue en memoire du processus, ce qui ne resisterait
  pas a une mise a l'echelle horizontale. \\
Absence de terminaison TLS &
  La passerelle publie en clair. Le chiffrement devra etre ajoute avant
  toute mise en ligne. \\
Hote unique &
  Le deploiement repose sur Docker Compose, sur une seule machine. Il n'y
  a ni orchestrateur de grappe ni redondance. \\
\bottomrule
\end{tabularx}
\end{table}

Ces ecarts sont le prix d'un projet mene en \stageDuree{} mois par une
seule personne. Ils sont documentes dans les diagrammes d'architecture
eux-memes, sous forme de notes, afin qu'un lecteur ne les prenne pas pour
des oublis.

% =====================================================================
\section*{Conclusion}
\addcontentsline{toc}{section}{Conclusion}

Ce chapitre a decrit les moyens et l'architecture. La plateforme repose
sur quatre niveaux logiques, deployes en vingt conteneurs sur un hote
unique. Chaque service metier applique le decoupage en couches de la
Clean Architecture, dont la regle de dependance protege le domaine des
choix techniques. Deux usages du courtier de messages coexistent
volontairement, pour ne pas imposer aux services Python un format concu
pour .NET. Quatre ecarts par rapport au patron de reference ont enfin ete
identifies et justifies.

Les six chapitres suivants presentent la realisation, sprint par sprint.
Le premier d'entre eux traite du socle sur lequel tout le reste
s'appuie : les identites et les acces.
